\subsubsection{নথির গোপনীয়তায় একই ধারণার প্রয়োগ}
ডেটাবেজ থেকে গ্রাহকের তালিকা রপ্তানি করে Word নথিতে রাখলেও সেই তথ্যের গোপনীয়তা রক্ষা করতে হয়। Windows-এর সমর্থিত ডেস্কটপ Word-এ File, Info, Protect Document ও Encrypt with Password নির্বাচন করে পাসওয়ার্ড দেওয়া যায়। পাসওয়ার্ড নিশ্চিত করার পরে নথিটি সংরক্ষণ করলে ফাইল খোলার সুরক্ষা প্রয়োগ হয়।

শুধু ফাইলের নাম বা এক্সটেনশন বদলানো, XML (Extensible Markup Language) হিসেবে সংরক্ষণ করা অথবা সম্পাদনার ওপর বিধিনিষেধ দেওয়া একই ধরনের এনক্রিপশন নিশ্চিত করে না। Word নথি এনক্রিপ্ট করা হলেও উৎস ডেটাবেজের অনুমতি ও ব্যাকআপের ব্যবস্থা প্রয়োজন থাকে।

